\chapter{Проектирование интеллектуальной системы классификации электронной почты}\label{ch:ch2}

В этой главе описываются данные,архитектура приложения,
порядок обработки писем, участие пользователя в~разметке
и~схема хранения состояния.

\section{Цель и границы проектирования}\label{sec:ch2/goal}

Система должна получать текстовые~документы из~почтового ящика,
преобразовывать их в~единый текстовый вид, строить векторные
представления, находить группы похожих документов, передавать
представителей групп на~ручную проверку и~после этого классифицировать новые
письма с~оценкой уверенности.

Проектируемая система не~заменяет пользователя при~задании смысла категорий.
Автоматические методы подготавливают кандидаты в~классы и~сокращают объём
ручной работы, а~пользователь подтверждает итоговые метки. Такой подход
соответствует постановке из~раздела~\ref{sec:ch1/dev}: если готовой разметки
мало или~нет, сначала нужно найти группы похожих документов, затем получить
проверенные метки и~только после этого строить классификатор.

\section{Выбор электронной почты как набора данных}\label{sec:ch2/dataset}

Электронная почта выбрана как предметная область, потому что она хорошо
представляет задачу классификации текстовых документов. Письмо содержит
естественный текст, тему, технические заголовки, дату, отправителя и~часто
относится к~понятной пользователю теме: учёба, работа, покупки, уведомления,
документы, личная переписка. При этом поток писем постоянно пополняется, поэтому
ручная сортировка со~временем требует всё больше повторяющихся действий.

Почтовый ящик также удобен технически. Большинство почтовых сервисов
поддерживает получение сообщений через стандартный почтовый клиент, а~каждое
сообщение имеет устойчивый идентификатор в~рамках папки. Это позволяет строить
синхронизацию как инкрементальный процесс: система загружает новые письма,
сохраняет их в~локальном хранилище и~не~обрабатывает уже известные сообщения
повторно.

Также электронная почта важна ещё и~как реалистичный источник
шумных данных. В~архиве встречаются короткие уведомления, цепочки ответов,
HTML-разметка, подписи, цитаты предыдущих сообщений, ссылки, даты, адреса и
служебные строки.

\section{Требования к системе и учёт ограничений}\label{sec:ch2/requirements}

Система должна работать с~личным почтовым архивом. Поэтому главным ограничением
становится локальная обработка: полный текст письма не~передаётся во~внешние
сервисы классификации. Внешние зависимости допускаются только как локально
запущенные компоненты или~как загруженные модели, выполняемые на~машине
пользователя. Это требование снижает риск раскрытия содержания писем и
сохраняет контроль над~данными.

Второе ограничение связано с~разметкой. В~почтовом ящике обычно нет готового
набора классов, пригодного для обучения модели. Папки почтового клиента могут
быть неполными, часть писем остаётся во~входящих, а~пользовательские правила
часто задаются техническими признаками, а~не~смыслом текста. Поэтому система
должна поддерживать режим начального формирования классов через кластеризацию
и~ручную проверку представителей.

Третье ограничение связано с~ошибками автоматической классификации. Письмо
может не~относиться ни~к~одной известной категории или~быть слишком коротким
для~надёжного решения. По~этой причине система должна использовать отказ
от~классификации, описанный в~разделе~\ref{sec/multitask}. Если максимальная
уверенность ниже заданного порога, письмо получает служебную метку
\texttt{other} и~может быть рассмотрено при~следующем обновлении.

\begin{table}[htbp]
    \centering
    \caption{Основные требования к проектируемой системе}\label{tab:ch2/requirements}
    \begin{tabular}{@{}p{0.27\textwidth}p{0.63\textwidth}@{}}
        \toprule
        Требование & Содержание \\
        \midrule
        Локальность & Тексты писем, нормализованные представления и метки хранятся
        в локальной базе данных. \\
        Воспроизводимость & Каждый этап сохраняет входные и выходные данные, чтобы
        обработку можно было повторить при изменении модели или параметров. \\
        Работа без полной разметки & Первичные классы формируются через
        кластеризацию похожих писем и ручную проверку представителей. \\
        Отказ от слабого решения & При низкой уверенности классификатор
        возвращает \texttt{other}, а не назначает случайную категорию. \\
        Обновляемость & Новые и шумовые письма накапливаются в буфере,
        после чего проходят повторную кластеризацию и разметку. \\
        \bottomrule
    \end{tabular}
\end{table}
\FloatBarrier

\section{Общая архитектура приложения}\label{sec:ch2/architecture}

Архитектура строится как последовательный конвейер обработки. Каждый этап
получает данные из~хранилища, выполняет одно преобразование и~записывает
результат обратно. Такой подход упрощает повторный запуск отдельных частей:
например, можно заменить модель эмбеддингов и~пересчитать только векторные
представления, не~загружая письма заново.

В~системе выделяются следующие уровни: получение данных, нормализация,
векторизация, кластеризация, разметка, классификация, оценка качества и
хранение состояния. Компоненты не~должны напрямую обращаться к~внутренним
структурам друг друга. Обмен идёт через записи в~PostgreSQL и~через явно
заданные идентификаторы писем, кластеров, меток и~версий классификатора.

Связь этапов с~общим хранилищем показана на~рисунке~\ref{fig:ch2/architecture}.

\begin{figure}[htbp]
    \centering
    \resizebox{\textwidth}{!}{%
        \begin{tikzpicture}[
        thick,
        >=stealth,
        node distance=6mm and 7mm,
        mod/.style={
                draw,
                rounded corners=2pt,
                align=center,
                minimum height=12mm,
                minimum width=24mm,
                inner sep=3pt,
                font=\small,
            },
        ctrl/.style={
                draw,
                rounded corners=2pt,
                align=center,
                minimum height=9mm,
                minimum width=40mm,
                inner sep=3pt,
                font=\small,
                fill=black!4,
            },
        store/.style={
                draw,
                rounded corners=2pt,
                align=center,
                minimum height=11mm,
                inner sep=3pt,
                font=\small,
                fill=black!4,
            },
    ]
    \node[ctrl]                                            (orch)  {orchestration\\(управление процессом)};
    \node[ctrl, below=of orch]                             (queue) {queueing\\очередь \texttt{pipeline\_tasks}};

    \node[mod, below=12mm of queue.south, anchor=north, xshift=-52mm] (email) {email\\синхронизация,\\нормализация};
    \node[mod, right=of email]                             (clu)   {clustering\\эмбеддинги,\\HDBSCAN / KMeans};
    \node[mod, right=of clu]                               (ann)   {annotation\\Argilla:\\экспорт / импорт};
    \node[mod, right=of ann]                               (clf)   {classifier\\прототипы,\\порог};
    \node[mod, right=of clf]                               (eval)  {evaluation\\метрики,\\дрейф};

    \node[store, below=12mm of email.south west, anchor=north west, minimum width=147mm]
                                                           (store) {storage "--- PostgreSQL\\(письма, нормализованные тексты, эмбеддинги, кластеры, метки, прототипы)};

    \draw[->] (orch) -- (queue);
    \draw[->] (email) -- (clu);
    \draw[->] (clu)   -- (ann);
    \draw[->] (ann)   -- (clf);
    \draw[->] (clf)   -- (eval);

    \draw[->] (queue.south) -| (email.north);
    \draw[->] (queue.south) -| (clf.north);

    \draw[<->, dashed] (email) -- (email |- store.north);
    \draw[<->, dashed] (clu)   -- (clu   |- store.north);
    \draw[<->, dashed] (ann)   -- (ann   |- store.north);
    \draw[<->, dashed] (clf)   -- (clf   |- store.north);
    \draw[<->, dashed] (eval)  -- (eval  |- store.north);
\end{tikzpicture}
%
    }
    \caption{Общая схема компонентов и хранилища проектируемой системы}\label{fig:ch2/architecture}
\end{figure}
\FloatBarrier

Прикладная часть реализуется на~Python, потому что для него доступны
библиотеки обработки текста и~машинного обучения. Для построения эмбеддингов
используется семейство моделей Sentence-BERT и~современные многоязычные
эмбеддеры~\autocite{Reimers2019SBERT,Reimers2020MultilingualSBERT,Chen2024BGEM3}.
Для кластеризации применяется HDBSCAN~\autocite{Campello2013HDBSCAN,McInnes2017HDBSCAN}, для расчёта
метрик и~вспомогательных алгоритмов "--- scikit-learn~\autocite{Pedregosa2011sklearn}.
Основная логика системы, инфраструктурные сервисы, такие как PostgreSQL и~сервис
ручной разметки,запускаются в~контейнерах. Контейнеры изолируют среду выполнения
программы и снижают зависимость результата от~локальной настройки машины.

\begin{table}[htbp]
    \centering
    \caption{Компоненты проектируемого приложения}\label{tab:ch2/components}
    \begin{tabular}{@{}p{0.27\textwidth}p{0.63\textwidth}@{}}
        \toprule
        Компонент & Назначение \\
        \midrule
        Синхронизация & Получает письма из почтового ящика и сохраняет сырой
        вариант сообщения. \\
        Нормализация & Извлекает полезный текст, удаляет технический шум и
        приводит письмо к единому виду. \\
        Векторизация & Строит эмбеддинг письма с помощью локально запущенной
        модели. \\
        Кластеризация & Группирует похожие письма и выделяет шумовые объекты. \\
        Разметка & Передаёт представителей кластеров пользователю и принимает
        проверенные категории. \\
        Классификация & Назначает категорию новым письмам или возвращает
        \texttt{other} при низкой уверенности. \\
        Хранилище & Сохраняет письма, эмбеддинги, кластеры, метки, версии
        классификатора и результаты обработки. \\
        \bottomrule
    \end{tabular}
\end{table}
\FloatBarrier

\section{Основной цикл программы}\label{sec:ch2/pipeline}

Основной цикл системы состоит из~двух режимов: начального построения
классификатора и~регулярного обновления. В~первом режиме система обрабатывает
уже имеющиеся письма в~почтовом ящике и~получает первую размеченную выборку.
Во~втором режиме она принимает новые письма, классифицирует надёжные случаи
и~накапливает спорные сообщения для~следующего витка разметки.

Автоматический режим работы показан на~рисунке~\ref{fig:ch2/pipeline}. После
полной синхронизации, нормализации и~построения эмбеддингов система проверяет
наличие начальных кластеров, накопление \texttt{other}, завершённые запуски
кластеризации и~готовую разметку. После классификации очередного пакета писем
без класса цикл возвращается к~проверке \texttt{other}.

\begin{figure}[htbp]
    \centering
    \resizebox{\textwidth}{!}{%
        \begin{tikzpicture}[
        thick,
        >=stealth,
        stage/.style={
                draw,
                rounded corners=2pt,
                align=center,
                text width=45mm,
                minimum height=10mm,
                inner sep=3pt,
                font=\scriptsize,
            },
        check/.style={
                stage,
                fill=black!4,
            },
        action/.style={
                stage,
                fill=black!8,
            },
        flow/.style={->, thick},
        note/.style={font=\scriptsize},
    ]

    \node[stage]  (sync)       at (0mm,    0mm) {Синхронизация\\всех писем};
    \node[stage]  (norm)       at (0mm,  -16mm) {Нормализация\\всех писем};
    \node[stage]  (embed)      at (0mm,  -32mm) {Построение\\эмбеддингов};
    \node[check]  (initcheck)  at (0mm,  -52mm) {Начальные\\кластеры есть?};
    \node[action] (initrun)    at (66mm, -52mm) {Первоначальная кластеризация писем};
    \node[check]  (othercheck) at (0mm,  -76mm) {В~\texttt{other}\\больше N писем?};
    \node[action] (otherrun)   at (66mm, -76mm) {Кластеризация\\писем \texttt{other}};
    \node[stage]  (export)     at (66mm,-100mm) {Экспорт\\кластеризации\\на разметку};
    \node[check]  (anncheck)   at (66mm,-124mm) {Есть завершённые\\неимпортированные\\размеченные данные?};
    \node[action] (importfit)  at (66mm,-148mm) {Импорт разметки\\и запуск\\\texttt{fit-prototypes}};
    \node[stage]  (classify)   at (0mm, -148mm) {Классификация\\K штук\\писем без класса};

    \draw[flow] (sync) -- (norm);
    \draw[flow] (norm) -- (embed);
    \draw[flow] (embed) -- (initcheck);

    \draw[flow] (initcheck.south) -- node[note, left] {да} (othercheck.north);
    \draw[flow] (initcheck.east) -- node[note, above] {нет} (initrun.west);
    \draw[flow] (initrun.east) -- ++(8mm, 0mm) |- (export.east);

    \draw[flow] (othercheck.south) -- node[note, left, pos=0.15] {нет} (classify.north);
    \draw[flow] (othercheck.east) -- node[note, above] {да} (otherrun.west);
    \draw[flow] (otherrun.south) -- (export.north);

    \draw[flow] (export) -- (anncheck);
    \draw[flow] (anncheck.east) .. controls +(18mm, 6mm) and +(18mm, -6mm) ..
        node[note, right, pos=0.5] {нет: ожидание} (anncheck.south east);
    \draw[flow] (anncheck.south) -- node[note, right] {да} (importfit.north);
    \draw[flow] (importfit.west) -- (classify.east);

    \draw[flow] (classify.west) -- ++(-28mm, 0mm) |- node[note, left, pos=0.25] {к шагу 5} (othercheck.west);
\end{tikzpicture}
%
    }
    \caption{Автоматический цикл обработки писем и обновления классификатора}\label{fig:ch2/pipeline}
\end{figure}
\FloatBarrier

Начальный цикл выполняется в~следующем порядке:

\begin{enumerate}
    \item загрузить письма из~почтового ящика и~сохранить их в~сыром виде;
    \item извлечь текст письма и~построить нормализованное представление;
    \item вычислить эмбеддинги нормализованных текстов;
    \item разделить эмбеддинги на~кластеры и~шум;
    \item выбрать представителей кластеров для~ручной проверки;
    \item получить пользовательские метки и~распространить их на~согласованные
        кластеры;
    \item построить классификатор по~размеченным письмам;
    \item зафиксировать версию классификатора и~протокол оценки.
\end{enumerate}

Ожидаемый результат цикла "--- набор писем с~проверенными метками, служебная
группа \texttt{other} для~неоднозначных сообщений, сохранённые эмбеддинги и
версия классификатора. После этого система может обрабатывать другие письма без
полной повторной разметки ранее загруженных писем.

\section{Синхронизация почтового ящика}\label{sec:ch2/sync}

Синхронизация отвечает за~получение сообщений из~почтового ящика и~защиту
от~повторной обработки. Для каждого письма система сохраняет технический
идентификатор, имя папки, дату получения, тему, адрес отправителя и~сырое
MIME-сообщение. Сырая версия нужна для~воспроизводимости: если изменится
алгоритм нормализации, письмо можно обработать повторно без нового обращения
к~почтовому серверу.

Процесс синхронизации должен быть идемпотентным. Повторный запуск не~создаёт
дубликаты, а~только добавляет письма, которых ещё нет в~локальной базе. Для
этого у~записи хранится внешний идентификатор сообщения и~папка, из~которой оно
получено. Если почтовый клиент переместил письмо между папками, система может
либо сохранить новую связь, либо считать письмо уже известным по~глобальному
идентификатору сообщения. Выбор конкретного правила относится к~реализации, но
проект должен предусматривать устойчивый ключ для~поиска дубликатов.

На~этапе синхронизации система не~принимает решений о~категории. Она только
сохраняет исходные данные и~ставит письмо в~очередь последующей обработки.
Такое разделение снижает последствия сетевых ошибок: сбой подключения к
почтовому серверу не~затрагивает уже нормализованные письма, эмбеддинги и
результаты классификации.

\section{Нормализация электронных писем}\label{sec:ch2/normalization}

Нормализация переводит письмо из~почтового формата в~текст, пригодный для
построения эмбеддинга. Система извлекает тему и~тело письма, выбирает
текстовую часть сообщения, удаляет HTML-разметку и~приводит пробельные символы
к~единому виду. Если письмо содержит только вложение или~короткую служебную
строку, оно помечается как непригодное для~основной классификации.

После извлечения текста система заменяет переменные элементы устойчивыми
токенами. Ссылки, адреса электронной почты, даты, время и~числа часто не
задают тему письма сами по~себе, но~увеличивают разнообразие строк. Замена
таких элементов на~\texttt{LINK}, \texttt{EMAIL}, \texttt{DATE},
\texttt{TIME} и~\texttt{NUM} уменьшает шум и~оставляет смысловую часть текста.
Этот этап развивает требования раздела~\ref{sec:ch1/preprocess}: модель должна
получать сопоставимые тексты, а~не~случайный набор технических строк.

Результатом нормализации является запись с~очищенным текстом и~оценкой
пригодности. Система хранит причину исключения, если письмо оказалось пустым,
слишком коротким или~техническим. Это позволяет не~терять такие сообщения, но
не~использовать их для~кластеризации и~проверки классификатора.

\section{Выбор модели и построение эмбеддингов}\label{sec:ch2/embeddings}

Для классификации писем важна смысловая близость, а~не~только совпадение
ключевых слов. Два письма могут относиться к~одной категории, хотя в~них
используются разные формулировки. Поэтому в~качестве основного представления
выбираются эмбеддинги: каждому нормализованному письму соответствует вектор
фиксированной размерности, по~которому можно вычислять близость документов.

Модель эмбеддингов должна удовлетворять трём условиям. Во-первых, она должна
работать локально, чтобы тексты не~передавались внешнему API. Во-вторых, она
должна поддерживать русский и~английский языки, поскольку в~почтовом архиве
могут встречаться оба языка. В-третьих, она должна давать представления,
подходящие для~сравнения текстов по~косинусной близости. Этим требованиям
соответствуют многоязычные модели семейства Sentence-BERT и~BGE-M3
\autocite{Reimers2019SBERT,Reimers2020MultilingualSBERT,Chen2024BGEM3}.

Выбор BGE-M3 дополнительно подтверждается прикладным сравнением открытых
моделей эмбеддингов. В~статье рассматривались девять open-source
моделей, которые запускались локально и~сравнивались по~десяти вопросам:
учитывались попадание релевантных документов в~первые пять результатов,
оценка качества ответа и~время обработки~\autocite{Bugra2025HabrEmbeddings}.
В~русскоязычном тесте модель \texttt{BAAI/bge-m3} получила среднюю оценку
3,50, как и~\texttt{multilingual-e5-large-instruct}, но~работала быстрее
(18,03~с против 21,47~с). В~смешанном тесте с~русским языком и~SQL лучший
результат показала \texttt{jina-embeddings-v3}, но~она обрабатывала запросы
дольше, чем \texttt{BAAI/bge-m3}. Более лёгкая \texttt{e5-small-v2} работала
быстрее, но~дала меньшую оценку качества. Для данной работы это означает, что
важен не~максимальный результат в~одном отдельном тесте, а~устойчивый баланс:
модель должна работать с~русским текстом, поддерживать локальный запуск
и~давать представления, пригодные для~сравнения по~косинусной близости. Поэтому
выбрана \texttt{BAAI/bge-m3}: она является открытой многоязычной моделью
и~сохраняет сбалансированное соотношение качества и~скорости.

Построенные эмбеддинги используются повторно на~нескольких этапах. HDBSCAN
получает их как вход для~поиска плотных групп, классификатор строит прототипы
классов по~эмбеддингам размеченных писем, а~модуль обновления сравнивает новые
сообщения с~уже известными классами. Поэтому эмбеддинг сохраняется в~хранилище
вместе с~версией модели. Если модель меняется, система должна пересчитать
векторы и~создать новую версию последующих результатов.

\section{Кластеризация и подбор параметров}\label{sec:ch2/clustering}

Кластеризация нужна на~этапе, когда у~системы ещё нет достаточной разметки.
Алгоритм получает эмбеддинги писем и~ищет группы, внутри которых сообщения
похожи друг на~друга сильнее, чем на~остальной корпус. Эти группы становятся
кандидатами в~пользовательские категории или~в~части более крупных категорий.

Для этой задачи выбран метод HDBSCAN. В~отличие от~K-Means,
он не~требует заранее задавать число кластеров. Это важно для~почтового архива:
пользователь заранее не~знает, сколько устойчивых тем есть в~уже накопленной
выборке писем. HDBSCAN также выделяет шумовые точки, что соответствует служебной
категории \texttt{other}: часть писем не~нужно принудительно относить к
какому-либо кластеру.

Различие между принудительным разбиением и~выделением шума показано на
рисунке~\ref{fig:ch2/clustering-outliers}.

\begin{figure}[htbp]
    \centering
    \resizebox{\textwidth}{!}{%
        \begin{tikzpicture}[
        thick,
        >=stealth,
        every node/.style={font=\scriptsize},
        cluster1/.style={circle, draw=black, fill=black!10, inner sep=1pt, minimum size=2.6mm},
        cluster2/.style={circle, draw=black, fill=black!30, inner sep=1pt, minimum size=2.6mm},
        cluster3/.style={circle, draw=black, fill=black!60, text=white, inner sep=1pt, minimum size=2.6mm},
        noise/.style={circle, draw=black, fill=white, inner sep=1pt, minimum size=2.4mm, dashed},
        panel/.style={draw, rounded corners=2pt, minimum width=56mm, minimum height=42mm},
    ]

    \begin{scope}[local bounding box=left]
        \node[panel] (Lframe) at (0,0) {};
        \node[anchor=south, font=\scriptsize, yshift=1.2mm] at (Lframe.north) {а) k-means: все точки приписаны};
        % три кластера + «выбросы», но все покрашены
        \foreach \x/\y in {-1.6/0.7, -1.2/1.1, -1.8/0.4, -1.4/0.2, -1.1/0.9}
            \node[cluster1] at (\x,\y) {};
        \foreach \x/\y in {0.6/0.5, 1.0/0.9, 0.4/1.0, 0.9/0.3, 1.2/0.6}
            \node[cluster2] at (\x,\y) {};
        \foreach \x/\y in {-0.4/-1.0, -0.1/-1.4, 0.2/-1.1, -0.3/-0.7}
            \node[cluster3] at (\x,\y) {};
        % «истинные выбросы», тоже покрашенные (искусственно отнесённые)
        \foreach \x/\y/\c in {1.7/-1.3/cluster3, -1.8/-1.2/cluster1, 1.8/1.4/cluster2}
            \node[\c] at (\x,\y) {};
    \end{scope}

    \begin{scope}[xshift=72mm, local bounding box=right]
        \node[panel] (Rframe) at (0,0) {};
        \node[anchor=south, font=\scriptsize, yshift=1.2mm] at (Rframe.north) {б) HDBSCAN: часть точек "--- шум};
        % те же три кластера
        \foreach \x/\y in {-1.6/0.7, -1.2/1.1, -1.8/0.4, -1.4/0.2, -1.1/0.9}
            \node[cluster1] at (\x,\y) {};
        \foreach \x/\y in {0.6/0.5, 1.0/0.9, 0.4/1.0, 0.9/0.3, 1.2/0.6}
            \node[cluster2] at (\x,\y) {};
        \foreach \x/\y in {-0.4/-1.0, -0.1/-1.4, 0.2/-1.1, -0.3/-0.7}
            \node[cluster3] at (\x,\y) {};
        % выбросы — отдельный стиль (шум)
        \foreach \x/\y in {1.7/-1.3, -1.8/-1.2, 1.8/1.4}
            \node[noise] at (\x,\y) {};
    \end{scope}
\end{tikzpicture}
%
    }
    \caption{Выделение шумовых точек при плотностной кластеризации}\label{fig:ch2/clustering-outliers}
\end{figure}
\FloatBarrier

Качество кластеризации зависит от~параметров алгоритма. Ручной выбор плохо
переносится между корпусами, поэтому проект предусматривает оптимизацию параметров
с помощью случайного подбора и закрепления лучшего результата по итогам оптимизации.
Целевая функция должна учитывать внутренние метрики кластеризации из
раздела~\ref{sec:ch1/metrics}, прежде всего DBCV~\autocite{Moulavi2014DBCV}, а
также практические ограничения: допустимую долю шума, минимальный размер
кластера и~разумный диапазон числа найденных групп. Такой подбор не~заменяет
ручную проверку, но~помогает выбрать конфигурацию, которая даёт пригодные
кандидаты для~разметки.

\section{Ручная проверка и назначение меток}\label{sec:ch2/annotation}

Кластеризация находит похожие письма, но~не~называет пользовательские категории.
Смысл метки задаёт человек. Чтобы не~размечать весь архив, система выбирает
несколько представителей каждого кластера: ближайшие к~центру группы письма,
письма с~высокой уверенностью принадлежности или~примеры, покрывающие разные
части кластера. Эти представители передаются в~интерфейс ручной проверки.

Пользователь назначает каждому представителю одну из~категорий или~служебную
метку \texttt{other}. Если представители одного кластера получили согласованную
метку, система распространяет её на~остальные письма этого кластера. Если
представители получили разные метки, кластер считается неоднородным. В~таком
случае система не~использует его целиком для~обучения: письма остаются
неразмеченными, попадают в~\texttt{other} или~требуют дополнительного
разделения.

Такой порядок сохраняет контроль пользователя над~содержанием классов. Модель
ускоряет подготовку выборки, но~не~создаёт категории без проверки. Для
последующего анализа система хранит источник метки: ручная разметка
представителя, автоматическое распространение по~кластеру или~служебное
решение из-за шума.

\section{Классификатор и цикл обновления}\label{sec:ch2/classifier}

После появления размеченного набора система строит классификатор поверх имеющихся
эмбеддингов. В~данной работе выбран прототипный подход, описанный в
разделе~\ref{sec:ch1/clf}. Для каждого класса вычисляются
прототипы по~эмбеддингам обучающих писем. Новое письмо нормализуется,
кодируется той~же моделью эмбеддингов и~сравнивается с~прототипами классов.
Класс ближайшего прототипа становится кандидатом на~ответ.

Поток применения классификатора показан на~рисунке~\ref{fig:ch2/classifier}.

\begin{figure}[htbp]
    \centering
    \resizebox{\textwidth}{!}{%
        \begin{tikzpicture}[
        thick,
        >=stealth,
        node distance=8mm and 10mm,
        block/.style={
                draw,
                rounded corners=2pt,
                align=center,
                minimum height=10mm,
                minimum width=28mm,
                inner sep=4pt,
                font=\small,
            },
        decision/.style={
                draw,
                diamond,
                aspect=2,
                align=center,
                inner sep=2pt,
                font=\small,
            },
        leaf/.style={
                draw,
                rounded corners=2pt,
                align=center,
                minimum height=9mm,
                minimum width=24mm,
                inner sep=4pt,
                font=\small,
                fill=black!4,
            },
    ]
    \node[block]                                    (msg)    {Сообщение \( x \in X \)\\(MIME, заголовки, тело)};
    \node[block, right=of msg]                      (norm)   {Нормализация\\\texttt{normalized\_text}};
    \node[block, right=of norm]                     (emb)    {Векторное\\представление \( \varphi(x) \)};
    \node[block, below=18mm of emb]                 (clf)    {Модель оценки\\уверенности\\\( g_\Theta \colon \mathbb{R}^{d} \to \mathbb{R}^{K} \)};
    \node[decision, left=of clf]                    (thr)    {\( \max_k s_k \geqslant \tau \)?};
    \node[leaf, left=of thr, yshift=10mm]           (label)  {Класс \( y_k \)\\\(+\)~оценка};
    \node[leaf, left=of thr, yshift=-10mm]          (other)  {\texttt{other}\\(буфер обратной связи)};

    \draw[->] (msg)  -- (norm);
    \draw[->] (norm) -- (emb);
    \draw[->] (emb)  -- (clf);
    \draw[->] (clf)  -- (thr);
    \draw[->] (thr.west) |- node[above, font=\scriptsize, pos=0.75] {да} (label.east);
    \draw[->] (thr.west) |- node[below, font=\scriptsize, pos=0.75] {нет} (other.east);
\end{tikzpicture}
%
    }
    \caption{Применение прототипного классификатора с порогом отказа}\label{fig:ch2/classifier}
\end{figure}
\FloatBarrier

Решающее правило включает порог уверенности. Если максимальная близость
недостаточна, система возвращает \texttt{other}. Это снижает число ошибочных
автоматических решений и~оставляет пользователю контроль над~необычными
письмами. При оценке качества нужно учитывать не~только точность на~письмах,
получивших категорию, но~и~покрытие, то~есть долю писем, для~которых система
не~отказалась от~ответа.

Дообучение системы строится вокруг накопления низкоуверенных писем.
Когда классификатор возвращает \texttt{other}, письмо сохраняется
вместе с~эмбеддингом и~оценками близости. После накопления достаточного числа
таких сообщений система запускает кластеризацию только этих писем.
Новые устойчивые группы отправляются на ручную проверку, после
чего классификатор получает новых кандидатов для обучения системы.

\section{Хранение данных в PostgreSQL}\label{sec:ch2/storage}

Хранилище должно фиксировать состояние каждого этапа обработки. Для этого
используется PostgreSQL: реляционная схема хорошо подходит для~писем,
нормализованных текстов, меток, версий классификатора и~связей между ними.
Большие векторные представления можно хранить как массивы, двоичные значения
или~внешние файлы с~ссылкой из~таблицы. Конкретный способ выбирается при
реализации, но~проект должен сохранять связь эмбеддинга с~письмом.

\begin{table}[htbp]
    \centering
    \caption{Основные сущности хранилища}\label{tab:ch2/storage}
    \begin{tabular}{@{}p{0.27\textwidth}p{0.63\textwidth}@{}}
        \toprule
        Сущность & Назначение \\
        \midrule
        Сырые письма & Записи \texttt{raw\_emails}: внешние идентификаторы,
        папки, даты и технические заголовки. \\
        Нормализованные письма & Записи \texttt{normalized\_emails}: очищенный
        текст, признаки качества и причина исключения из обработки. \\
        Эмбеддинги & Векторные представления писем и версия модели эмбеддингов. \\
        Кластеры & Номер запуска кластеризации, параметры алгоритма, метка
        кластера и признак шума. \\
        Метки & Пользовательские категории, источник метки и связь с письмом
        или кластером. \\
        Версии классификатора & Использованные данные, порог уверенности и
        служебные параметры версии. \\
        Предсказания & Результаты применения классификатора, уверенность ответа
        и факт отказа. \\
        \bottomrule
    \end{tabular}
\end{table}
\FloatBarrier

\section{Выводы по главе}\label{sec:ch2/concl}

Во~второй главе спроектирована локальная интеллектуальная система для
классификации электронной почты. Определены требования к~локальной обработке,
работе без полной разметки, отказу от~слабых решений и~обновлению модели при
поступлении новых писем.

Предложенная архитектура строится как последовательный конвейер: синхронизация
почтового ящика, нормализация сообщений, построение эмбеддингов, кластеризация,
ручная проверка, построение классификатора и~цикл обновления через буфер
шумовых объектов. Для хранения состояния выбрана реляционная схема PostgreSQL,
которая связывает исходные письма, промежуточные результаты, метки, версии
классификатора и~предсказания.
